\documentclass[11pt,a4paper]{article}
\usepackage{xeCJK}
\usepackage{fontspec}
\usepackage{xcolor}
\usepackage{geometry}
\usepackage[protrusion=true]{microtype}
\geometry{margin=1.8cm,top=1.9cm}
\linespread{1.25}

\setmainfont{Baskerville}
\setCJKmainfont{Songti SC}

\definecolor{tipblue}{RGB}{0,95,160}
\definecolor{rulegray}{RGB}{190,190,190}

\setlength{\parindent}{0pt}
\setlength{\parskip}{0.35em}
\newcommand{\num}[1]{\makebox[1.9em][r]{#1.}\hspace{0.3em}}
\newcommand{\wline}{\noindent\rule{\linewidth}{0.4pt}}
% 中文题干：宋体粗体；提示词：蓝色小字；英文句子：Baskerville 正体
\newcommand{\prompt}[1]{\emph{\textcolor{tipblue}{\small（#1）}}}

\begin{document}
\pagestyle{plain}
\begin{center}
{\LARGE\bfseries 上海高考高分翻译背诵 \quad 95}\\[0.25em]
{\large 24一模　2026-09-25 \quad\ 中英对照（背诵 · 答案）}\\[0.35em]
{\footnotesize 来源：微信公众号「发现之旅 速来记单词」· 2026-09-25 · 赵文瀚}
\end{center}

\vspace{0.4em}
\noindent\begin{tabular}{@{}p{0.33\textwidth}@{}p{0.34\textwidth}@{}p{0.33\textwidth}@{}}
姓名：\underline{\hspace{2.0cm}} & 日期：\underline{\hspace{2.0cm}} & 得分：\underline{\hspace{1.4cm}}\,/\enspace 4 \\
\end{tabular}

\vspace{0.5em}
\noindent{\small 中 $\to$ 英默写卷的参考答案；亦可直接用于背诵与回译自查。}

\vspace{0.8em}
\num{1}\textbf{如果不好好准备，周五的演讲可能会变得一塌糊涂。}\prompt{preparation}

\vspace{0.15em}

\noindent\hangindent=2.2em\hangafter=1\hspace{2.2em}The speech on Friday might go south without adequate preparation.

\num{2}\textbf{市民们呼吁废纸回收再利用,以减少对原材料的消耗。}\prompt{call for}

\vspace{0.15em}

\noindent\hangindent=2.2em\hangafter=1\hspace{2.2em}Citizens call for the recycling of waste paper to reduce the consumption of raw materials.

\num{3}\textbf{电影里出现了许多主人公穿越沙漠的场景，象征着一种自我发现和成长的过程。}\prompt{there}

\vspace{0.15em}

\noindent\hangindent=2.2em\hangafter=1\hspace{2.2em}There are many scenes in the film where the leading character crosses the desert, symbolizing a process of self-discovery and growth.

\num{4}\textbf{艺术博物馆位于中国著名的文化城市杭州，在那里游客们可以欣赏风格多样的艺术作品，仿佛置身于艺术的海洋之中。}\prompt{where}

\vspace{0.15em}

\noindent\hangindent=2.2em\hangafter=1\hspace{2.2em}Located in Hangzhou, a renowned cultural city in China, the art museum is a place where tourists can appreciate diverse styles of works as if immersed in an ocean of art.

\end{document}
